\documentclass[aip,jcp,reprint,noshowkeys,superscriptaddress]{revtex4-1}
\usepackage{graphicx,dcolumn,bm,xcolor,microtype,multirow,amscd,amsmath,amssymb,amsfonts,physics,longtable,wrapfig}
\usepackage[version=4]{mhchem}

\usepackage[utf8]{inputenc}
\usepackage[T1]{fontenc}
\usepackage{txfonts}

\usepackage[
    colorlinks,
    linkcolor={red!50!black},
    citecolor={red!70!black},
    urlcolor={red!80!black}
	]{hyperref}
\urlstyle{same}

\newcommand{\ie}{\textit{i.e.}}
\newcommand{\eg}{\textit{e.g.}}
\newcommand{\alert}[1]{\textcolor{red}{#1}}
\usepackage[normalem]{ulem}
\newcommand{\titou}[1]{\textcolor{red}{#1}}
\newcommand{\trashPFL}[1]{\textcolor{\red}{\sout{#1}}}
\newcommand{\PFL}[1]{\titou{(\underline{\bf PFL}: #1)}}

\newcommand{\cre}[1]{\Hat{a}_{#1}^{\dag}}
\newcommand{\ani}[1]{\Hat{a}_{#1}}
\newcommand{\ca}[2]{\Hat{a}^{#1}_{#2}}

\newcommand{\mc}{\multicolumn}
\newcommand{\fnm}{\footnotemark}
\newcommand{\fnt}{\footnotetext}
\newcommand{\tabc}[1]{\multicolumn{1}{c}{#1}}
\newcommand{\QP}{\textsc{quantum package}}
\newcommand{\T}[1]{#1^{\intercal}}

% coordinates
\newcommand{\br}{\boldsymbol{r}}
\newcommand{\bx}{\boldsymbol{x}}
\newcommand{\bI}{\boldsymbol{1}}
\newcommand{\cK}{\mathcal{K}}

% orbital energies
\newcommand{\eps}{\epsilon}
\newcommand{\irr}{\text{irr}}

% shortcuts for greek letters
\newcommand{\si}{\sigma}
\newcommand{\la}{\lambda}
\newcommand{\ii}{\mathrm{i}}
\newcommand{\up}{\uparrow}
\newcommand{\dw}{\downarrow}
\newcommand{\upup}{\up\up}
\newcommand{\updw}{\up\dw}
\newcommand{\dwup}{\dw\up}
\newcommand{\dwdw}{\dw\downarrow}

\newcommand{\h}{\text{1h}}
\newcommand{\p}{\text{1p}}
\newcommand{\hp}{\text{1h+1p}}
\newcommand{\hhp}{\text{2h1p}}
\newcommand{\pph}{\text{2p1h}}
\newcommand{\ph}{\text{ph}}
\newcommand{\eh}{\text{eh}}
\newcommand{\he}{\text{he}}
\newcommand{\ee}{\text{ee}}
\newcommand{\teh}{\overline{\text{eh}}}
\newcommand{\pp}{\text{pp}}
\newcommand{\hh}{\text{hh}}
\newcommand{\xc}{\text{xc}}
\newcommand{\Hx}{\text{Hx}}

% addresses
\newcommand{\LCPQ}{Laboratoire de Chimie et Physique Quantiques (UMR 5626), Universit\'e de Toulouse, CNRS, UPS, France}

\begin{document}

\title{Multiresolvent Representation of the Parquet Vertex through Third Order}

\author{Pierre-Fran\c{c}ois \surname{Loos}}
	\email{loos@irsamc.ups-tlse.fr}
	\affiliation{\LCPQ}

\begin{abstract}
A general four-point vertex depends on three independent frequencies and cannot, in general, be represented by the single resolvent familiar from linear-response theory or conventional algebraic-diagrammatic construction (ADC).
In these notes, we develop a channel-resolved multiresolvent representation of the parquet vertex through third order in the interaction.
The electron-hole ($\eh$) and particle-particle ($\pp$) reducible vertices $\Phi^\eh$ and $\Phi^\pp$ are chosen as the two primitive objects because their central propagation occurs naturally in the $N$- and $(N\pm2)$-electron sectors, respectively, while the transverse electron-hole ($\teh$) contribution follows from crossing symmetry.
In each primitive channel, the reducible vertex is represented as two frequency-dependent external vertices connected by a central resolvent.
The direct $\eh$ irreducible vertex, and therefore the corresponding four-point Hedin kernel, is reconstructed from $\Gamma^\eh=\Lambda+\Phi^{\teh}+\Phi^\pp$.
At second order, integrating over one fermionic frequency of the exact three-frequency kernel with the corresponding noninteracting pair propagator produces two-frequency external vertices whose poles belong exclusively to the $(N\pm1)$-electron sectors.
Integrating over the remaining fermionic frequency recovers the known one-frequency Bethe--Salpeter equation (BSE) kernels associated with second-order Green's-function theory (GF2) in both the $\eh$ and $\pp$ channels.
At third order, the direct $\eh$ reducible vertex is generated by second-order corrections to the two external vertices together with a first-order correction to the central neutral resolvent.
The analogous $\pp$ construction has the same external-vertex structure but a central resolvent associated with the $(N\pm2)$-electron sectors.
The resulting formulation provides a compact recipe for reconstructing the full parquet vertex and the Hedin kernel from ordinary resolvents associated with the $N$-, $(N\pm1)$-, and $(N\pm2)$-electron sectors.
\end{abstract}

\maketitle

\tableofcontents

%%%%%%%%%%%%%%%%%%%%%%%%%
\section{Introduction}
\label{sec:introduction}
%%%%%%%%%%%%%%%%%%%%%%%%%

A general four-point quantity contains substantially more analytic
structure than a one-particle Green's function or a linear-response
propagator.
For a time-independent Hamiltonian, time-translation invariance removes
one overall time variable, leaving three independent time differences or,
equivalently, three independent frequencies.
The full two-particle vertex can therefore be written as
$F(\nu,\nu',\omega)$.

This three-frequency dependence makes a direct extension of conventional
spectral or algebraic-diagrammatic construction (ADC) non-trivial.
A one-particle self-energy can be represented by resolvents in the
$(N\pm1)$-electron sectors, while a polarization propagator is naturally
expressed through a neutral $N$-electron resolvent.
For a four-point vertex, several propagation sectors and several independent
frequencies coexist.
The objective is therefore not to introduce a single formal resolvent
carrying all three frequencies, but to organize the vertex into ordinary
resolvents associated with well-defined intermediate-state spaces.

Parquet theory provides the natural channel decomposition for this purpose.
We denote the electron-hole, transverse electron-hole, and particle-particle
channels by $\eh$, $\teh$, and $\pp$, respectively.
The full vertex is decomposed as\cite{DeDominicis_1964b,Bickers_2004,MarieLoos2025}
\begin{equation}
	F
	=
	\Lambda
	+
	\Phi^\eh
	+
	\Phi^{\teh}
	+
	\Phi^\pp
	\label{eq:parquet_decomposition}
\end{equation}
where $\Lambda$ is fully two-particle irreducible.
The direct $\eh$ reducible vertex $\Phi^\eh$ is associated with neutral
$N$-electron propagation, whereas $\Phi^\pp$ is associated with the
$(N\pm2)$-electron sectors.
The transverse contribution $\Phi^{\teh}$ follows from $\Phi^\eh$ by
fermionic crossing.
Consequently, only two primitive multiresolvent constructions are required:
$\Phi^\eh$ and $\Phi^\pp$.

In each primitive channel, the channel-reducible four-point vertex is
represented as a left external vertex, a channel-specific central resolvent,
and a right external vertex
\begin{equation}
	\text{left external vertex}
	\;\times\;
	\text{central resolvent}
	\;\times\;
	\text{right external vertex}
	\label{eq:architecture}
\end{equation}
The central resolvent carries the natural bosonic transfer of the channel,
while the two external vertices retain the remaining fermionic-frequency
dependence.

This construction also provides a direct route to the four-point Hedin
kernel.
The vertex irreducible in the direct $\eh$ channel is
\begin{equation}
	\Gamma^\eh
	=
	\Lambda
	+
	\Phi^{\teh}
	+
	\Phi^\pp
	\label{eq:Gamma_eh_parquet}
\end{equation}
At the exact level, this is the same four-point kernel that enters Hedin's
Bethe--Salpeter equation (BSE) through the functional derivative
$\Xi=\fdv*{\Sigma}{G}$, once a common index convention is adopted\cite{Hedin1965,MarieAmmarLoos2024}.
Thus, constructing $\Phi^\eh$ and $\Phi^\pp$ also yields a multiresolvent
representation of the direct $\eh$ Hedin kernel.

The presentation is deliberately separated by perturbative order.
We first establish a complete second-order construction: channel frequency
structure, primitive reducible vertices, one-sided frequency integrations,
external-vertex pole spaces, recovery of the second-order Green's-function
(GF2) BSE kernels, and parquet consistency.
A dedicated recipe then summarizes the complete result through second order.
Only after this point do we introduce the additional ingredients required at
third order: second-order corrections to the two external vertices and the
first-order correction to the central resolvent.

Throughout the perturbative discussion, we use a Hartree--Fock (HF) reference
and a M\o ller--Plesset (MP) partition.
For compactness, all expressions are written for real-valued orbitals.
Left and right external vertices are kept distinct whenever this distinction
matters, and all causal prescriptions are denoted by $0^+$ because integer
multiples of the infinitesimal regulator carry no physical significance.

\section{Frequency conventions and channel organization}
\label{sec:frequencies}
%%%%%%%%%%%%%%%%%%%%%%%%%

The direct $\eh$ convention is used throughout these notes.
The four external fermionic line energies are
\begin{equation}
	\nu
	\qquad
	\nu+\omega
	\qquad
	\nu'
	\qquad
	\nu'+\omega
	\label{eq:eh_frequencies}
\end{equation}
The natural bosonic transfer frequencies of the three parquet channels are
\begin{align}
	\Omega_\eh
	&=
	\omega
	\label{eq:Omega_eh}
	\\
	\Omega_{\teh}
	&=
	\nu'-\nu
	\label{eq:Omega_ehbar}
	\\
	\Omega_\pp
	&=
	\omega+\nu+\nu'
	\label{eq:Omega_pp}
\end{align}

For the transverse $\eh$ channel, exchanging the third and fourth external
legs gives the frequency map
\begin{equation}
	\left(
	\nu_{\teh},
	\nu'_{\teh},
	\Omega_{\teh}
	\right)
	=
	\left(
	\nu,
	\nu+\omega,
	\nu'-\nu
	\right)
	\label{eq:crossed_frequency_map}
\end{equation}
The natural $\pp$ variables are related to the direct $\eh$ convention by
\begin{equation}
	\left(
	\nu_\pp,
	\nu'_\pp,
	\Omega_\pp
	\right)
	=
	\left(
	-\omega-\nu',
	-\nu',
	\omega+\nu+\nu'
	\right)
	\label{eq:pp_frequency_map}
\end{equation}

These three combinations determine the characteristic frequency structure of
the parquet vertex.
At the lowest non-vanishing order, each reducible vertex depends only on the
bosonic transfer natural to its channel
\begin{align}
	\Phi^{\eh,(2)}
	&\equiv
	\Phi^{\eh,(2)}(\omega)
	\\
	\Phi^{\teh,(2)}
	&\equiv
	\Phi^{\teh,(2)}(\nu'-\nu)
	\\
	\Phi^{\pp,(2)}
	&\equiv
	\Phi^{\pp,(2)}(\nu+\nu'+\omega)
	\label{eq:Phi_frequency_second_order}
\end{align}
Beyond second order, the primitive reducible vertices are genuine
three-frequency objects.
The multiresolvent representation introduced below assigns the natural
bosonic transfer to the central resolvent and the two remaining fermionic
frequencies to the external vertices.

The exact channel-irreducible vertices satisfy
\begin{align}
	\Gamma^\eh
	&=
	\Lambda
	+
	\Phi^{\teh}
	+
	\Phi^\pp
	\label{eq:Gamma_eh}
	\\
	\Gamma^\pp
	&=
	\Lambda
	+
	\Phi^\eh
	+
	\Phi^{\teh}
	\label{eq:Gamma_pp}
\end{align}
For a two-body interaction, the first correction to the fully irreducible
vertex beyond the bare interaction occurs at fourth order, so that
\begin{equation}
	\Lambda
	=
	\Lambda^{(1)}
	+
	O(v^4)
	\label{eq:Lambda_order}
\end{equation}
The second-order channel-irreducible vertices are therefore
\begin{align}
	\Gamma^{\eh,(2)}
	&=
	\Phi^{\teh,(2)}
	+
	\Phi^{\pp,(2)}
	\label{eq:Gamma_eh_second}
	\\
	\Gamma^{\pp,(2)}
	&=
	\Phi^{\eh,(2)}
	+
	\Phi^{\teh,(2)}
	\label{eq:Gamma_pp_second}
\end{align}

\section{Target multiresolvent representation}
\label{sec:target_representation}
%%%%%%%%%%%%%%%%%%%%%%%%%

Having fixed the channel frequencies, we can state the algebraic form that we
seek before deriving any perturbative ingredient.

In conventional ADC for the self-energy, one starts from a diagonal spectral
representation in the exact $(N\pm1)$-electron eigenstate basis and uses the
unitary freedom within the intermediate-state space to obtain a non-diagonal
ADC representation.
For a general four-point vertex, no analogous single diagonal representation
exists because three independent frequencies cannot be carried by one
resolvent.
There is therefore no direct analogue of the usual ADC transformation from a
diagonal three-frequency representation to a non-diagonal one.

Our strategy is instead channel resolved.
For each primitive reducible vertex, the natural bosonic transfer is carried
by a central resolvent and the remaining fermionic-frequency dependences by
left and right external vertices.
The target forms are
\begin{align}
	\ii\Phi^\eh(\nu,\nu',\omega)
	&=
	U_L^\eh(\nu,\omega)
	R^\eh(\omega)
	U_R^\eh(\nu',\omega)
	\label{eq:target_eh}
	\\
	\ii\Phi_P^\pp(\nu_\pp,\nu'_\pp,\Omega_\pp)
	&=
	U_L^\pp(\nu_\pp,\Omega_\pp)
	R^\pp(\Omega_\pp)
	U_R^\pp(\nu'_\pp,\Omega_\pp)
	\label{eq:target_pp}
\end{align}
with channel-specific central resolvents
\begin{align}
	R^\eh(\omega)
	&=
	\left[
	\omega\eta^\eh-M^\eh
	\right]^{-1}
	\label{eq:target_R_eh}
	\\
	R^\pp(\Omega_\pp)
	&=
	\left[
	\Omega_\pp\eta^\pp-M^\pp
	\right]^{-1}
	\label{eq:target_R_pp}
\end{align}
Here $\eta^r$ is the metric of channel $r$, with eigenvalues $+1$ and $-1$
for its forward and backward sectors, respectively.
The symbols $U_L^r$ and $U_R^r$ denote external vertices and should not be
confused with the unitary transformation used in conventional ADC.

The corresponding particle-number structure is
\begin{align}
	\Phi^\eh
	&:\quad
	N\pm1
	\longrightarrow
	N
	\longrightarrow
	N\pm1
	\label{eq:target_sector_eh}
	\\
	\Phi^\pp
	&:\quad
	N\pm1
	\longrightarrow
	N\pm2
	\longrightarrow
	N\pm1
	\label{eq:target_sector_pp}
\end{align}
Thus, the central resolvent distinguishes neutral and pair propagation,
whereas the external vertices are associated with $(N\pm1)$-electron
propagation.
At second order, their non-trivial poles are the $\hhp$ and $\pph$ poles of
the $(N-1)$- and $(N+1)$-electron sectors.

The zeroth-order central resolvent is written generically as
\begin{equation}
	R^{r,(0)}(\Omega_r)
	=
	\left[
	\Omega_r\eta^r-M^{r,(0)}
	\right]^{-1}
	\label{eq:target_R_zero}
\end{equation}
At second order, the proposed representation must reduce to
\begin{equation}
	\ii\Phi^{r,(2)}
	=
	U_L^{r,(1)}
	R^{r,(0)}
	U_R^{r,(1)}
	\label{eq:target_second_order}
\end{equation}
The first part of these notes establishes this second-order representation
completely and validates the associated one-sided frequency integrations.
The extension of the same ansatz to third order is introduced only after the
second-order recipe in Sec.~\ref{sec:recipe_second}.

\section{Three-frequency BSE}
\label{sec:three_frequency_bse}
%%%%%%%%%%%%%%%%%%%%%%%%%

The exact $\eh$ correlation function obeys the four-point BSE
\begin{equation}
	L
	=
	L_0
	+
	L_0\Gamma^\eh L
	\label{eq:BSE_symbolic}
\end{equation}
where products imply integrations over the internal space-spin-time
variables.
Because a four-point quantity depends on three independent time differences,
the frequency-space BSE contains two internal fermionic-frequency
integrations
\begin{multline}
	L(\nu,\nu',\omega)
	=
	L_0(\nu,\nu',\omega)
	\\
	+
	\frac{1}{(2\pi)^2}
	\int d\bar\nu d\bar\nu'
	L_0(\nu,\bar\nu,\omega)
	\Gamma^\eh(\bar\nu,\bar\nu',\omega)
	L(\bar\nu',\nu',\omega)
	\label{eq:BSE_three_frequency}
\end{multline}

For the noninteracting reference, the $\eh$ pair propagator is diagonal in
the fermionic frequency
\begin{equation}
	L_0(\nu,\nu',\omega)
	=
	2\pi\delta(\nu-\nu')
	G_0(\nu+\omega)G_0(\nu)
	\label{eq:L0_three_frequency}
\end{equation}
The ordering in Eq.~\eqref{eq:L0_three_frequency} fixes our convention: a
forward pair $i\to a$ carries the one-particle frequencies $\nu$ and
$\nu+\omega$ through $G_{0,i}$ and $G_{0,a}$, respectively.
For occupied and virtual orbitals,
\begin{align}
	G_{0,i}(\nu)
	&=
	\frac{1}{\nu-\eps_i-\ii\eta}
	&
	G_{0,a}(\nu)
	&=
	\frac{1}{\nu-\eps_a+\ii\eta}
	\label{eq:G0}
\end{align}

%%%%%%%%%%%%%%%%%%%%%%%%%
\subsection{One-sided factorization of the pair propagator}
%%%%%%%%%%%%%%%%%%%%%%%%%


The diagonal structure of the noninteracting three-frequency propagator
allows its frequency dependence to be separated explicitly.
In the forward and backward eh sectors,
\begin{equation}
	L_{0,ia}^{\sigma}(\nu,\nu',\omega)
	=
	2\pi\delta(\nu-\nu')
	Q_{ia}^{\sigma}(\nu,\omega)
	(L_0^\sigma)_{ia}(\omega)
	\qquad
	\sigma=\pm
	\label{eq:L0_factorization}
\end{equation}
Thus, the three-frequency propagator factorizes into a normalized
two-frequency factor $Q^\sigma(\nu,\omega)$ and the usual one-frequency
eh propagator $(L_0^\sigma)(\omega)$.

We therefore introduce the normalized one-sided factors
\begin{align}
	Q^+_{ia}(\nu,\omega)
	&=
	-\ii
	\left[
	G_{0,i}(\nu)-G_{0,a}(\nu+\omega)
	\right]
	\label{eq:Q_plus}
	\\
	Q^-_{ia}(\nu,\omega)
	&=
	+\ii
	\left[
	G_{0,a}(\nu)-G_{0,i}(\nu+\omega)
	\right]
	\label{eq:Q_minus}
\end{align}
Partial fractioning of the two one-particle propagators then gives the
factorizations
\begin{align}
	G_{0,a}(\nu+\omega)G_{0,i}(\nu)
	&=
	Q^+_{ia}(\nu,\omega)
	(L_0^+)_{ia}(\omega)
	\label{eq:L0_factorization_plus}
	\\
	G_{0,i}(\nu+\omega)G_{0,a}(\nu)
	&=
	Q^-_{ia}(\nu,\omega)
	(L_0^-)_{ia}(\omega)
	\label{eq:L0_factorization_minus}
\end{align}
Thus, $L_0^\pm$ carries the bosonic transfer $\omega$, whereas $Q^\pm$
retains the fermionic frequency that will be integrated on one side of the
BSE.
With the infinitesimal time shifts implied in the time-ordered frequency
integrations,
\begin{equation}
	\frac{1}{2\pi}
	\int d\nu\,
	Q^\pm_{ia}(\nu,\omega)
	=
	1
	\label{eq:Q_normalization}
\end{equation}
and pair reversal gives
\begin{equation}
	Q^-_{ia}(\nu,\omega)
	=
	Q^+_{ia}(\nu+\omega,-\omega)
	\label{eq:Q_reversal}
\end{equation}

%%%%%%%%%%%%%%%%%%%%%%%%%
\subsection{One-sided BSE vertices}
%%%%%%%%%%%%%%%%%%%%%%%%%

Equations~\eqref{eq:L0_factorization_plus} and
\eqref{eq:L0_factorization_minus} explain how the external vertices arise.
Instead of integrating both fermionic frequencies of the three-frequency
kernel, we integrate only one side of the BSE and keep the other fermionic
frequency explicit.
Using the convention that the factor $\ii$ is absorbed into the
channel-irreducible vertex, a right-sided BSE factor has the form
\begin{equation}
	\begin{split}
	&\frac{1}{2\pi}
	\int d\nu\,
	Q^\sigma_{ia}(\nu,\omega)
	(L_0^\sigma)_{ia}(\omega)
	\left[
	\ii\Gamma^\eh
	\right]^{\sigma\sigma'}_{ia,jb}
	(\nu,\nu',\omega)
	\\
	&\qquad=
	(L_0^\sigma)_{ia}(\omega)
	U_R^{\sigma\sigma'}{}_{ia,jb}(\nu',\omega)
	\end{split}
	\label{eq:BSE_right_factorization}
\end{equation}
which defines
\begin{equation}
	U_R^{\sigma\sigma'}{}_{ia,jb}(\nu',\omega)
	=
	\frac{1}{2\pi}
	\int d\nu\,
	Q^\sigma_{ia}(\nu,\omega)
	\left[
	\ii\Gamma^\eh
	\right]^{\sigma\sigma'}_{ia,jb}
	(\nu,\nu',\omega)
	\label{eq:U_R_def}
\end{equation}
Likewise, the corresponding left-sided factorization is
\begin{equation}
	\begin{split}
	&\frac{1}{2\pi}
	\int d\nu'\,
	\left[
	\ii\Gamma^\eh
	\right]^{\sigma\sigma'}_{ia,jb}
	(\nu,\nu',\omega)
	Q^{\sigma'}_{jb}(\nu',\omega)
	(L_0^{\sigma'})_{jb}(\omega)
	\\
	&\qquad=
	U_L^{\sigma\sigma'}{}_{ia,jb}(\nu,\omega)
	(L_0^{\sigma'})_{jb}(\omega)
	\end{split}
	\label{eq:BSE_left_factorization}
\end{equation}
with
\begin{equation}
	U_L^{\sigma\sigma'}{}_{ia,jb}(\nu,\omega)
	=
	\frac{1}{2\pi}
	\int d\nu'\,
	\left[
	\ii\Gamma^\eh
	\right]^{\sigma\sigma'}_{ia,jb}
	(\nu,\nu',\omega)
	Q^{\sigma'}_{jb}(\nu',\omega)
	\label{eq:U_L_def}
\end{equation}
Here $\sigma,\sigma'=+,-$ label the forward and backward pair sectors.
The objects in Eqs.~\eqref{eq:U_R_def} and \eqref{eq:U_L_def} are therefore
not independent ansatzes: they are the factors that remain after the
one-frequency propagator $L_0^\sigma(\omega)$ has been extracted from one
side of the three-frequency BSE.
They are the pair-space form of the frequency-dependent external vertices:
one fermionic frequency has been integrated out, while the other fermionic
frequency and the bosonic transfer remain explicit.
The same one-sided construction, with arbitrary external orbital pairs,
produces the $U_L^r$ and $U_R^r$ entering the multiresolvent ansatz of
Sec.~\ref{sec:target_representation}.
Moreover, the normalization in Eq.~\eqref{eq:Q_normalization} guarantees
that an instantaneous first-order kernel is left unchanged by the one-sided
integration, which leads directly to the bare first-order external vertices
used below.

The frequency reduction can therefore be summarized as
\begin{equation}
	\Gamma^\eh(\nu,\nu',\omega)
	\longrightarrow
	\left\{
	U_L(\nu,\omega),
	U_R(\nu',\omega)
	\right\}
	\longrightarrow
	\widetilde\Gamma^\eh(\omega)
	\label{eq:contraction_hierarchy}
\end{equation}
where the first arrow denotes a one-sided frequency integration and the
second denotes integration over the remaining fermionic frequency.

%%%%%%%%%%%%%%%%%%%%%%%%%

\section{Second-order primitive vertices}
\label{sec:second_order_primitives}
%%%%%%%%%%%%%%%%%%%%%%%%%

We now construct the second-order $\eh$ irreducible vertex from
the two channels that are irreducible with respect to direct $\eh$
propagation.
Antisymmetrized two-electron integrals are denoted as
\begin{equation}
	V_{pqrs}
	=
	\mel{pq}{}{rs}
	\label{eq:V_def}
\end{equation}

%%%%%%%%%%%%%%%%%%%%%%%%%
\subsection{Direct eh reducible vertex}
%%%%%%%%%%%%%%%%%%%%%%%%%

At first order, the $\eh$ irreducible interaction is the bare
antisymmetrized interaction
\begin{equation}
	\ii\Gamma^{\eh,(1)}_{pqrs}
	=
	V_{pqrs}
	\label{eq:Gamma_eh_first}
\end{equation}
The second-order direct $\eh$ reducible vertex follows from the
zeroth-order neutral excitations $i\to a$
\begin{equation}
	\begin{split}
	\ii\Phi^{\eh,(2)}_{pqrs}(\omega)
	={}&
	\sum_{ia}
	\frac{
	V_{pa\,ri}
	V_{qi\,sa}
	}
	{\omega-(\eps_a-\eps_i)+\ii0^+}
	\\
	&-
	\sum_{ia}
	\frac{
	V_{pi\,ra}
	V_{qa\,si}
	}
	{\omega+(\eps_a-\eps_i)-\ii0^+}
	\end{split}
	\label{eq:Phi_eh_second}
\end{equation}
At this order the two end vertices are instantaneous and
$\Phi^{\eh,(2)}$ depends only on the direct $\eh$ transfer
frequency $\omega$.

%%%%%%%%%%%%%%%%%%%%%%%%%
\subsection{Transverse eh contribution from crossing}
%%%%%%%%%%%%%%%%%%%%%%%%%

Fermionic crossing gives
\begin{equation}
	\Phi^{\teh}_{pqrs}(\nu,\nu',\omega)
	=
	-
	\Phi^\eh_{pqsr}
	\left(
	\nu,
	\nu+\omega,
	\nu'-\nu
	\right)
	\label{eq:Phi_ehbar_crossing}
\end{equation}
Using Eq.~\eqref{eq:Phi_eh_second} therefore yields
\begin{equation}
	\begin{split}
	\ii\Phi^{\teh,(2)}_{pqrs}(\nu,\nu',\omega)
	={}&
	-
	\sum_{kc}
	\frac{
	V_{pc\,sk}
	V_{qk\,rc}
	}
	{\nu'-\nu-(\eps_c-\eps_k)+\ii0^+}
	\\
	&+
	\sum_{kc}
	\frac{
	V_{pk\,sc}
	V_{qc\,rk}
	}
	{\nu'-\nu+(\eps_c-\eps_k)-\ii0^+}
	\end{split}
	\label{eq:Phi_ehbar_second}
\end{equation}
The characteristic transverse $\eh$ poles are therefore controlled
by $\nu'-\nu$.

%%%%%%%%%%%%%%%%%%%%%%%%%
\subsection{pp reducible vertex}
%%%%%%%%%%%%%%%%%%%%%%%%%

The $\pp$ channel contains $\pp$ and $\hh$
propagation.
Using unrestricted orbital sums, the second-order pair-reducible vertex
is
\begin{equation}
	\begin{split}
	\ii\Phi_P^{\pp,(2)}{}_{pqrs}(\Omega)
	={}&
	\frac{1}{2}
	\sum_{cd}
	\frac{
	V_{pq\,cd}
	V_{cd\,rs}
	}
	{\Omega-(\eps_c+\eps_d)+\ii0^+}
	\\
	&-
	\frac{1}{2}
	\sum_{kl}
	\frac{
	V_{pq\,kl}
	V_{kl\,rs}
	}
	{\Omega-(\eps_k+\eps_l)-\ii0^+}
	\end{split}
	\label{eq:Phi_pp_second_native}
\end{equation}
The factors $1/2$ remove the double counting associated with unrestricted
antisymmetric pair sums.
They are absent when the sums are restricted to unique pairs.

Transforming Eq.~\eqref{eq:Phi_pp_second_native} to the direct
$\eh$ convention with Eq.~\eqref{eq:pp_frequency_map} gives
\begin{equation}
	\begin{split}
	\ii\Phi^{\pp,(2)}_{pqrs}(\nu,\nu',\omega)
	={}&
	\frac{1}{2}
	\sum_{cd}
	\frac{
	V_{pq\,cd}
	V_{cd\,rs}
	}
	{\nu+\nu'+\omega-\eps_c-\eps_d+\ii0^+}
	\\
	&-
	\frac{1}{2}
	\sum_{kl}
	\frac{
	V_{pq\,kl}
	V_{kl\,rs}
	}
	{\nu+\nu'+\omega-\eps_k-\eps_l-\ii0^+}
	\end{split}
	\label{eq:Phi_pp_second_eh}
\end{equation}
The characteristic pair poles are therefore controlled by
$\nu+\nu'+\omega$.

%%%%%%%%%%%%%%%%%%%%%%%%%
\subsection{Second-order eh irreducible kernel}
%%%%%%%%%%%%%%%%%%%%%%%%%

Combining the transverse $\eh$ and $\pp$
contributions gives
\begin{equation}
	\ii\Gamma^{\eh,(2)}
	=
	\ii\Phi^{\teh,(2)}
	+
	\ii\Phi^{\pp,(2)}
	\label{eq:Gamma_eh_second_i}
\end{equation}
Explicitly,
\begin{equation}
	\begin{split}
	\ii\Gamma^{\eh,(2)}_{pqrs}(\nu,\nu',\omega)
	={}&
	-
	\sum_{kc}
	\frac{
	V_{pc\,sk}
	V_{qk\,rc}
	}
	{\nu'-\nu-(\eps_c-\eps_k)+\ii0^+}
	\\
	&+
	\sum_{kc}
	\frac{
	V_{pk\,sc}
	V_{qc\,rk}
	}
	{\nu'-\nu+(\eps_c-\eps_k)-\ii0^+}
	\\
	&+
	\frac{1}{2}
	\sum_{cd}
	\frac{
	V_{pq\,cd}
	V_{cd\,rs}
	}
	{\nu+\nu'+\omega-\eps_c-\eps_d+\ii0^+}
	\\
	&-
	\frac{1}{2}
	\sum_{kl}
	\frac{
	V_{pq\,kl}
	V_{kl\,rs}
	}
	{\nu+\nu'+\omega-\eps_k-\eps_l-\ii0^+}
	\end{split}
	\label{eq:Gamma_eh_second_explicit}
\end{equation}
Equation~\eqref{eq:Gamma_eh_second_explicit} is the starting point for
the one-sided frequency integration.

%%%%%%%%%%%%%%%%%%%%%%%%%
%%%%%%%%%%%%%%%%%%%%%%%%%
\section{Multiresolvent representation at second order}
\label{sec:second_order_multiresolvent}
%%%%%%%%%%%%%%%%%%%%%%%%%

The primitive second-order vertices can now be written explicitly in the
target form of Eq.~\eqref{eq:target_second_order}.
At this order, the external vertices are the bare interaction and the central
resolvent is noninteracting.

For the $\eh$ channel, the metric and zeroth-order neutral matrix are
\begin{align}
	\eta^\eh
	&=
	\begin{pmatrix}
	1&0
	\\
	0&-1
	\end{pmatrix}
	\label{eq:eta_eh}
	\\
	M^{\eh,(0)}
	&=
	\begin{pmatrix}
	\Delta\eps&0
	\\
	0&\Delta\eps
	\end{pmatrix}
	\label{eq:M_eh_zero}
\end{align}
with
\begin{equation}
	(\Delta\eps)_{ia,jb}
	=
	(\eps_a-\eps_i)\delta_{ij}\delta_{ab}
	\label{eq:Delta_eh}
\end{equation}
and
\begin{equation}
	R^{\eh,(0)}(\omega)
	=
	\left[
	\omega\eta^\eh-M^{\eh,(0)}
	\right]^{-1}
	\label{eq:R_eh_zero}
\end{equation}
The first-order left and right external vertices are
\begin{align}
	(U_L^{\eh,(1),+})_{pr,ia}
	&=
	V_{pa\,ri}
	&
	(U_R^{\eh,(1),+})_{ia,sq}
	&=
	V_{qi\,sa}
	\label{eq:U_eh_first_plus}
	\\
	(U_L^{\eh,(1),-})_{pr,ai}
	&=
	V_{pi\,ra}
	&
	(U_R^{\eh,(1),-})_{ai,sq}
	&=
	V_{qa\,si}
	\label{eq:U_eh_first_minus}
\end{align}
so that
\begin{equation}
	\ii\Phi^{\eh,(2)}
	=
	U_L^{\eh,(1)}
	R^{\eh,(0)}
	U_R^{\eh,(1)}
	\label{eq:Phi_eh_second_resolvent}
\end{equation}
which is identical to Eq.~\eqref{eq:Phi_eh_second}.

For the $\pp$ channel, we use a normalized antisymmetric pair basis with
unique particle pairs $a<b$ and hole pairs $i<j$.
The metric and zeroth-order pair matrix are
\begin{align}
	\eta^\pp
	&=
	\begin{pmatrix}
	1&0
	\\
	0&-1
	\end{pmatrix}
	\label{eq:eta_pp}
	\\
	M^{\pp,(0)}
	&=
	\begin{pmatrix}
	\eps_a+\eps_b&0
	\\
	0&-(\eps_i+\eps_j)
	\end{pmatrix}
	\label{eq:M_pp_zero}
\end{align}
with
\begin{equation}
	R^{\pp,(0)}(\Omega_\pp)
	=
	\left[
	\Omega_\pp\eta^\pp-M^{\pp,(0)}
	\right]^{-1}
	\label{eq:R_pp_zero}
\end{equation}
The first-order external vertices are
\begin{align}
	(U_L^{\pp,(1),+})_{pq,ab}
	&=
	V_{pq\,ab}
	&
	(U_R^{\pp,(1),+})_{ab,rs}
	&=
	V_{ab\,rs}
	\label{eq:U_pp_first_plus}
	\\
	(U_L^{\pp,(1),-})_{pq,ij}
	&=
	V_{pq\,ij}
	&
	(U_R^{\pp,(1),-})_{ij,rs}
	&=
	V_{ij\,rs}
	\label{eq:U_pp_first_minus}
\end{align}
Therefore
\begin{equation}
	\ii\Phi_P^{\pp,(2)}
	=
	U_L^{\pp,(1)}
	R^{\pp,(0)}
	U_R^{\pp,(1)}
	\label{eq:Phi_pp_second_resolvent}
\end{equation}
which reproduces Eq.~\eqref{eq:Phi_pp_second_native}.
The opposite signs of the $\pp$ and $\hh$ denominators are
encoded by the $\pm1$ eigenvalues of $\eta^\pp$.

Equations~\eqref{eq:Phi_eh_second_resolvent} and
\eqref{eq:Phi_pp_second_resolvent} establish the primitive multiresolvent
representation completely through second order.

\section{Pair-space external vertices at second order}
\label{sec:half_contraction_second}
%%%%%%%%%%%%%%%%%%%%%%%%%

The one-sided frequency integration removes one fermionic frequency while retaining
the other and the bosonic transfer frequency.
The resulting poles can be identified before considering the detailed
orbital numerators.
For a forward external pair $i\to a$, the relevant contour integrations
transform the transverse $\eh$ poles according to
\begin{equation}
	\begin{split}
	&\frac{1}{2\pi}
	\int d\nu
	Q^+_{ia}(\nu,\omega)
	\frac{1}
	{\nu'-\nu+(\eps_c-\eps_k)-\ii0^+}
	\\
	&\qquad=
	\frac{1}
	{\nu'-(\eps_i+\eps_k-\eps_c)-\ii0^+}
	\end{split}
	\label{eq:contour_ehbar_hhp}
\end{equation}
\begin{equation}
	\begin{split}
	&\frac{1}{2\pi}
	\int d\nu
	Q^+_{ia}(\nu,\omega)
	\frac{1}
	{\nu'-\nu-(\eps_c-\eps_k)+\ii0^+}
	\\
	&\qquad=
	\frac{1}
	{\nu'+\omega-(\eps_a+\eps_c-\eps_k)+\ii0^+}
	\end{split}
	\label{eq:contour_ehbar_pph}
\end{equation}
The $\pp$ poles give
\begin{equation}
	\begin{split}
	&\frac{1}{2\pi}
	\int d\nu
	Q^+_{ia}(\nu,\omega)
	\frac{1}
	{\nu+\nu'+\omega-\eps_c-\eps_d+\ii0^+}
	\\
	&\qquad=
	\frac{1}
	{\nu'+\omega-(\eps_c+\eps_d-\eps_i)+\ii0^+}
	\end{split}
	\label{eq:contour_pp_pph}
\end{equation}
\begin{equation}
	\begin{split}
	&\frac{1}{2\pi}
	\int d\nu
	Q^+_{ia}(\nu,\omega)
	\frac{1}
	{\nu+\nu'+\omega-\eps_k-\eps_l-\ii0^+}
	\\
	&\qquad=
	\frac{1}
	{\nu'-(\eps_k+\eps_l-\eps_a)-\ii0^+}
	\end{split}
	\label{eq:contour_pp_hhp}
\end{equation}

The signs of the corresponding contributions to the external vertices are
therefore inherited from the coefficients of
$\Gamma^{\eh,(2)}$ in Eq.~\eqref{eq:Gamma_eh_second_explicit}.
In particular, the negative $\pph$ transverse contribution and the
negative $\hhp$ pair contribution do not originate from the contour
integration itself.

The only energies appearing as poles of the external vertices are therefore
\begin{align}
	E^-_{ij a}
	&=
	\eps_i+\eps_j-\eps_a
	\label{eq:E_minus}
	\\
	E^+_{ab i}
	&=
	\eps_a+\eps_b-\eps_i
	\label{eq:E_plus}
\end{align}
which correspond to $\hhp$ states in the $N-1$ sector and $\pph$ states
in the $N+1$ sector.

%%%%%%%%%%%%%%%%%%%%%%%%%
\subsection{Forward-forward block}
%%%%%%%%%%%%%%%%%%%%%%%%%

The forward-forward block is identified as
\begin{equation}
	\left[
	\ii\Gamma^\eh
	\right]^{++}_{ia,jb}
	=
	\ii\Gamma^\eh_{aj\,ib}
	\label{eq:Gamma_pp_block}
\end{equation}
At first order,
\begin{equation}
	U^{++,(1)}_{ia,jb}
	=
	V_{aj\,ib}
	\label{eq:U_pp_first}
\end{equation}
At second order, inserting Eq.~\eqref{eq:Gamma_eh_second_explicit} in
Eq.~\eqref{eq:U_R_def} gives
\begin{equation}
	\begin{split}
	U^{++,(2)}_{ia,jb}(\nu,\omega)
	={}&
	\sum_{kc}
	\frac{
	V_{jc\,ik}
	V_{ka\,cb}
	}
	{\nu-(\eps_i+\eps_k-\eps_c)-\ii0^+}
	\\
	&-
	\sum_{kc}
	\frac{
	V_{jk\,ic}
	V_{ca\,kb}
	}
	{\nu+\omega-(\eps_a+\eps_c-\eps_k)+\ii0^+}
	\\
	&-
	\frac{1}{2}
	\sum_{kl}
	\frac{
	V_{aj\,kl}
	V_{lk\,bi}
	}
	{\nu-(\eps_k+\eps_l-\eps_a)-\ii0^+}
	\\
	&+
	\frac{1}{2}
	\sum_{cd}
	\frac{
	V_{aj\,cd}
	V_{dc\,bi}
	}
	{\nu+\omega-(\eps_c+\eps_d-\eps_i)+\ii0^+}
	\end{split}
	\label{eq:U_pp_second}
\end{equation}
The first two terms originate from the transverse $\eh$ channel,
while the last two originate from the $\hh$ and $\pp$
branches of $\Phi^\pp$.

%%%%%%%%%%%%%%%%%%%%%%%%%
\subsection{Forward-backward block}
%%%%%%%%%%%%%%%%%%%%%%%%%

The forward-backward block is
\begin{equation}
	\left[
	\ii\Gamma^\eh
	\right]^{+-}_{ia,jb}
	=
	\ii\Gamma^\eh_{ab\,ij}
	\label{eq:Gamma_pm_block}
\end{equation}
At first order,
\begin{equation}
	U^{+-,(1)}_{ia,jb}
	=
	V_{ab\,ij}
	\label{eq:U_pm_first}
\end{equation}
At second order,
\begin{equation}
	\begin{split}
	U^{+-,(2)}_{ia,jb}(\nu,\omega)
	={}&
	\sum_{kc}
	\frac{
	V_{bc\,ik}
	V_{ka\,cj}
	}
	{\nu-(\eps_i+\eps_k-\eps_c)-\ii0^+}
	\\
	&-
	\sum_{kc}
	\frac{
	V_{bk\,ic}
	V_{ca\,kj}
	}
	{\nu+\omega-(\eps_a+\eps_c-\eps_k)+\ii0^+}
	\\
	&-
	\frac{1}{2}
	\sum_{kl}
	\frac{
	V_{ab\,kl}
	V_{lk\,ij}
	}
	{\nu-(\eps_k+\eps_l-\eps_a)-\ii0^+}
	\\
	&+
	\frac{1}{2}
	\sum_{cd}
	\frac{
	V_{ab\,cd}
	V_{dc\,ij}
	}
	{\nu+\omega-(\eps_c+\eps_d-\eps_i)+\ii0^+}
	\end{split}
	\label{eq:U_pm_second}
\end{equation}
The same $\hhp$ and $\pph$ pole spaces therefore appear in both upper
pair-orientation blocks.

%%%%%%%%%%%%%%%%%%%%%%%%%
\subsection{Backward blocks}
%%%%%%%%%%%%%%%%%%%%%%%%%

For real orbitals, pair reversal relates the remaining two blocks to the
upper blocks
\begin{align}
	U^{--}_{ai,bj}(\nu,\omega)
	&=
	\left[
	U^{++}_{ia,jb}(\nu+\omega,-\omega)
	\right]^*
	\label{eq:U_mm_reversal}
	\\
	U^{-+}_{ai,jb}(\nu,\omega)
	&=
	\left[
	U^{+-}_{ia,jb}(\nu+\omega,-\omega)
	\right]^*
	\label{eq:U_mp_reversal}
\end{align}
The fermionic frequency carried by a given particle-number sector is therefore
controlled by the orientation of the external pair
\begin{equation}
	\begin{array}{c|cc}
	& \hhp & \pph
	\\
	\hline
	+ & \nu & \nu+\omega
	\\
	- & \nu+\omega & \nu
	\end{array}
	\label{eq:side_frequency_table}
\end{equation}
A derivation for complex orbitals requires the corresponding explicit
Hermitian and crossing relations and will not be needed below.

%%%%%%%%%%%%%%%%%%%%%%%%%
\section{Residue-tensor representation}
\label{sec:residue_tensor}
%%%%%%%%%%%%%%%%%%%%%%%%%

The second-order external vertices admit a compact spectral representation.
We first use ordered intermediate configurations
\begin{align}
	I^-
	&=
	(\mu\nu;\alpha)
	\in
	\hhp
	&
	E^-_{\mu\nu;\alpha}
	&=
	\eps_\mu+\eps_\nu-\eps_\alpha
	\label{eq:E_minus_general}
	\\
	I^+
	&=
	(\alpha\beta;\mu)
	\in
	\pph
	&
	E^+_{\alpha\beta;\mu}
	&=
	\eps_\alpha+\eps_\beta-\eps_\mu
	\label{eq:E_plus_general}
\end{align}
The ordered-tuple convention is used here only to keep the $1/2$ factors
of the pair channel explicit.
A transformation to normalized antisymmetric ADC intermediate states can
be performed afterwards.

For a forward pair, the general second-order structure is
\begin{equation}
	U^{\sigma\sigma',(2)}_{AB}(\nu,\omega)
	=
	\sum_{I^-}
	\frac{
	\mathcal R^{\sigma\sigma',-}_{AB;I^-}
	}
	{\nu-E^-_{I^-}-\ii0^+}
	+
	\sum_{I^+}
	\frac{
	\mathcal R^{\sigma\sigma',+}_{AB;I^+}
	}
	{\nu+\omega-E^+_{I^+}+\ii0^+}
	\label{eq:U_residue_general}
\end{equation}

For the forward-forward block, the $\hhp$ residue is
\begin{equation}
	\begin{split}
	\mathcal R^{++,-}_{ia,jb;\mu\nu;\alpha}
	={}&
	\delta_{\mu i}
	V_{j\alpha\,\mu\nu}
	V_{\nu a\,\alpha b}
	\\
	&-
	\frac{1}{2}
	\delta_{\alpha a}
	V_{aj\,\mu\nu}
	V_{\nu\mu\,bi}
	\end{split}
	\label{eq:R_pp_minus}
\end{equation}
and the $\pph$ residue is
\begin{equation}
	\begin{split}
	\mathcal R^{++,+}_{ia,jb;\alpha\beta;\mu}
	={}&
	-
	\delta_{\alpha a}
	V_{j\mu\,i\beta}
	V_{\beta a\,\mu b}
	\\
	&+
	\frac{1}{2}
	\delta_{\mu i}
	V_{aj\,\alpha\beta}
	V_{\beta\alpha\,bi}
	\end{split}
	\label{eq:R_pp_plus}
\end{equation}

For the forward-backward block, the $\hhp$ residue is
\begin{equation}
	\begin{split}
	\mathcal R^{+-,-}_{ia,jb;\mu\nu;\alpha}
	={}&
	\delta_{\mu i}
	V_{b\alpha\,\mu\nu}
	V_{\nu a\,\alpha j}
	\\
	&-
	\frac{1}{2}
	\delta_{\alpha a}
	V_{ab\,\mu\nu}
	V_{\nu\mu\,ij}
	\end{split}
	\label{eq:R_pm_minus}
\end{equation}
and the $\pph$ residue is
\begin{equation}
	\begin{split}
	\mathcal R^{+-,+}_{ia,jb;\alpha\beta;\mu}
	={}&
	-
	\delta_{\alpha a}
	V_{b\mu\,i\beta}
	V_{\beta a\,\mu j}
	\\
	&+
	\frac{1}{2}
	\delta_{\mu i}
	V_{ab\,\alpha\beta}
	V_{\beta\alpha\,ij}
	\end{split}
	\label{eq:R_pm_plus}
\end{equation}
The lower-block residues follow from pair reversal.

Equations~\eqref{eq:R_pp_minus}--\eqref{eq:R_pm_plus} reveal an important
feature of the multiresolvent construction.
The pole spaces of the external vertices are universally the $\hhp$ and $\pph$ sectors, but
the residues retain the topology of the parquet channel from which they
originate.
The transverse $\eh$ terms factor naturally in the transverse
pairing of the external indices, while the $\pp$ terms factor
naturally in the pair channel.
Consequently, a representation of the form
$U=CRD$ in the direct $\eh$ pair grouping is not generally the
minimal representation of the exact residue structure.
The safer second-order statement is the matrix-valued spectral form of
Eq.~\eqref{eq:U_residue_general}.

A simple vertex--resolvent--vertex factorization can still be recovered
channel by channel in the native pairing of each contribution or by
introducing a larger auxiliary space containing the required spectator
labels.
Neither construction is required at second order and the residue-tensor
form will be retained below.

%%%%%%%%%%%%%%%%%%%%%%%%%
\section{Second frequency integration: the eh GF2 kernel}
\label{sec:GF2_contraction}
\label{sec:second_contraction}
%%%%%%%%%%%%%%%%%%%%%%%%%

The remaining fermionic frequency can now be integrated with the
one-sided propagator of the second external pair.
This yields the BSE kernel associated with second-order Green's-function
theory (GF2).
For a forward closure, the elementary $\hhp$ and $\pph$ frequency integrations are
\begin{equation}
	\frac{1}{2\pi}
	\int d\nu
	Q^+_{jb}(\nu,\omega)
	\frac{1}
	{\nu-E^- -\ii0^+}
	=
	-\frac{1}
	{\omega-(\eps_b-E^-)+\ii0^+}
	\label{eq:second_contraction_minus}
\end{equation}
\begin{equation}
	\frac{1}{2\pi}
	\int d\nu
	Q^+_{jb}(\nu,\omega)
	\frac{1}
	{\nu+\omega-E^+ +\ii0^+}
	=
	\frac{1}
	{\omega-(E^+-\eps_j)+\ii0^+}
	\label{eq:second_contraction_plus}
\end{equation}
Applying these identities to Eq.~\eqref{eq:U_pp_second} gives the
second-order resonant block of the effective kernel
\begin{equation}
	\begin{split}
	(\widetilde\Xi_c^{(2)})^{jb}_{ia}(\omega)
	={}&
	-
	\sum_{kc}
	\frac{
	V_{jc\,ik}
	V_{ka\,cb}
	}
	{\omega-(\eps_b-\eps_i+\eps_c-\eps_k)}
	\\
	&-
	\sum_{kc}
	\frac{
	V_{jk\,ic}
	V_{ca\,kb}
	}
	{\omega-(\eps_a-\eps_j+\eps_c-\eps_k)}
	\\
	&+
	\frac{1}{2}
	\sum_{kl}
	\frac{
	V_{aj\,kl}
	V_{lk\,bi}
	}
	{\omega-(\eps_a+\eps_b-\eps_k-\eps_l)}
	\\
	&+
	\frac{1}{2}
	\sum_{cd}
	\frac{
	V_{aj\,cd}
	V_{dc\,bi}
	}
	{\omega-(\eps_c+\eps_d-\eps_i-\eps_j)}
	\end{split}
	\label{eq:GF2_A}
\end{equation}

Closing the forward-backward external vertex with a backward pair gives
\begin{equation}
	\begin{split}
	(\widetilde\Xi_c^{(2)})^{bj}_{ia}
	={}&
	-
	\sum_{kc}
	\frac{
	V_{bc\,ik}
	V_{ka\,cj}
	}
	{-(\eps_b-\eps_i+\eps_c-\eps_k)}
	\\
	&-
	\sum_{kc}
	\frac{
	V_{bk\,ic}
	V_{ca\,kj}
	}
	{-(\eps_a-\eps_j+\eps_c-\eps_k)}
	\\
	&+
	\frac{1}{2}
	\sum_{kl}
	\frac{
	V_{ab\,kl}
	V_{lk\,ij}
	}
	{-(\eps_a+\eps_b-\eps_k-\eps_l)}
	\\
	&+
	\frac{1}{2}
	\sum_{cd}
	\frac{
	V_{ab\,cd}
	V_{dc\,ij}
	}
	{-(\eps_c+\eps_d-\eps_i-\eps_j)}
	\end{split}
	\label{eq:GF2_B}
\end{equation}
The external transfer frequency cancels in the coupling block.
This is a consequence of closing the resolvents entering the external vertices with a backward
external pair rather than a new type of intermediate-state propagation.

The complete effective kernel therefore has the familiar Casida-like
structure
\begin{equation}
	\widetilde{\boldsymbol\Xi}^{(2)}(\omega)
	=
	\begin{pmatrix}
	A_c^{(2)}(\omega)
	&
	B_c^{(2)}
	\\
	B_c^{(2)\dagger}
	&
	A_c^{(2)\dagger}(-\omega)
	\end{pmatrix}
	\label{eq:GF2_matrix}
\end{equation}
Equations~\eqref{eq:GF2_A} and \eqref{eq:GF2_B} reproduce the effective
dynamic kernel obtained independently by differentiating the GF2
self-energy and performing the conventional two-frequency integration.
This agreement provides a stringent check of the channel frequency maps,
fermionic signs, and pair-channel factors.

The conventional effective kernel involves an additional approximation at
the level of the exact three-frequency BSE, where the
full correlation function appearing inside the frequency integration is replaced by
$L_0$.
The comparison made here is therefore with this $L_0$-projected dynamic
kernel rather than with the exact uncontracted three-frequency kernel.

%%%%%%%%%%%%%%%%%%%%%%%%%
%%%%%%%%%%%%%%%%%%%%%%%%%
\section{Second frequency integration: the pp GF2 kernel}
\label{sec:pp_gf2_check}
%%%%%%%%%%%%%%%%%%%%%%%%%

The $\pp$ construction admits an independent second-order check analogous
to the $\eh$ GF2 test of Sec.~\ref{sec:GF2_contraction}.
Starting from the second-order $\pp$-irreducible kernel in
Eq.~\eqref{eq:Gamma_pp_second}, integrating over one fermionic frequency with
the corresponding noninteracting pair propagator produces a two-frequency
pair-space external vertex $U^{\pp,(2)}$.
Contracting the remaining relative fermionic frequency must then recover
the one-frequency second-order pp kernel obtained independently from the
GF2 anomalous self-energy.
This comparison provides a stringent test because it probes simultaneously
the $\pph$ and $\hhp$ external-vertex poles, the antisymmetry of the pair space, and
the relative signs of the four orbital attachments.

Consider first an incoming two-particle pair $ab$.
After integrating over one fermionic frequency, the upper external vertex can be written as
\begin{equation}
\begin{split}
U_{ab,rs}^{\pp,(2)}(\bar\nu,\Omega)
={}&
\sum_{me}
\frac{
V_{rm\,ae}V_{se\,bm}
}
{\bar\nu-(\eps_a+\eps_e-\eps_m-\Omega/2)+\ii0^+}
\\
&+
\sum_{me}
\frac{
V_{re\,bm}V_{sm\,ae}
}
{\bar\nu-(\eps_m-\eps_a-\eps_e+\Omega/2)-\ii0^+}
\\
&-
\sum_{me}
\frac{
V_{re\,am}V_{sm\,be}
}
{\bar\nu-(\eps_m-\eps_e-\eps_b+\Omega/2)-\ii0^+}
\\
&-
\sum_{me}
\frac{
V_{rm\,be}V_{se\,am}
}
{\bar\nu-(\eps_b+\eps_e-\eps_m-\Omega/2)+\ii0^+}
\end{split}
\label{eq:pp_half_coupling_explicit}
\end{equation}
The two individual fermionic energies of the outgoing pair are
\begin{equation}
 z_1=\frac{\Omega}{2}+\bar\nu
 \qquad
 z_2=\frac{\Omega}{2}-\bar\nu
 \label{eq:pp_individual_energies}
\end{equation}
so the poles of Eq.~\eqref{eq:pp_half_coupling_explicit} occur at $\pph$
energies.
The two terms associated with $z_1$ and $z_2$ correspond to attaching the
particle-hole correction to either member of the antisymmetric pair.
The $\hh$ external vertex is obtained in the same manner and contains the
corresponding $\hhp$ energies.

Projecting the remaining pair onto a two-particle state $cd$ yields the
dynamic upper-left block of the effective pp kernel
\begin{equation}
\begin{split}
\ii\widetilde\Xi_{ab,cd}^{\pp,(2)}(\Omega)
={}&
\ii\sum_{me}
\frac{
V_{cm\,ae}V_{de\,bm}
}
{\Omega-(\eps_a+\eps_c+\eps_e-\eps_m)+\ii0^+}
\\
&-
\ii\sum_{me}
\frac{
V_{ce\,bm}V_{dm\,ae}
}
{\Omega-(\eps_a+\eps_d+\eps_e-\eps_m)+\ii0^+}
\\
&-
\ii\sum_{me}
\frac{
V_{cm\,be}V_{de\,am}
}
{\Omega-(\eps_b+\eps_c+\eps_e-\eps_m)+\ii0^+}
\\
&+
\ii\sum_{me}
\frac{
V_{ce\,am}V_{dm\,be}
}
{\Omega-(\eps_b+\eps_d+\eps_e-\eps_m)+\ii0^+}
\end{split}
\label{eq:pp_GF2_C}
\end{equation}
This is the correlation contribution to the dynamic $C$ block of the GF2
$\pp$ BSE matrix.
The four residues appear with the relative pattern $+,-,-,+$, exactly as
required by antisymmetry under interchange of either external pair.

Projecting instead onto a two-hole state $ij$ produces the coupling block
\begin{equation}
\begin{split}
\ii\widetilde\Xi_{ab,ij}^{\pp,(2)}
={}&
\ii\sum_{me}
\frac{
V_{ie\,bm}V_{jm\,ae}
}
{-(\eps_i-\eps_a+\eps_m-\eps_e)}
\\
&+
\ii\sum_{me}
\frac{
V_{im\,ae}V_{je\,bm}
}
{-(\eps_a-\eps_j+\eps_e-\eps_m)}
\\
&-
\ii\sum_{me}
\frac{
V_{ie\,am}V_{jm\,be}
}
{-(\eps_i-\eps_b+\eps_m-\eps_e)}
\\
&-
\ii\sum_{me}
\frac{
V_{im\,be}V_{je\,am}
}
{-(\eps_b-\eps_j+\eps_e-\eps_m)}
\end{split}
\label{eq:pp_GF2_B}
\end{equation}
The explicit dependence on the bosonic pair frequency cancels in this
block.
The static character of the $\ee$--$\hh$ coupling is therefore not an
additional approximation but follows directly from the analytic structure
of the partially integrated kernel.
This is the precise $\pp$ analogue of the static $B$ block in
the $\eh$ GF2 kernel.

Starting from the two-hole external vertex and integrating the remaining fermionic frequency with
a two-hole projector gives
\begin{equation}
\begin{split}
\ii\widetilde\Xi_{ij,kl}^{\pp,(2)}(-\Omega)
={}&
\ii\sum_{me}
\frac{
V_{ke\,im}V_{lm\,je}
}
{-\Omega-(-\eps_i-\eps_k+\eps_e-\eps_m)-\ii0^+}
\\
&-
\ii\sum_{me}
\frac{
V_{ke\,jm}V_{lm\,ie}
}
{-\Omega-(-\eps_j-\eps_k+\eps_e-\eps_m)-\ii0^+}
\\
&-
\ii\sum_{me}
\frac{
V_{km\,je}V_{le\,im}
}
{-\Omega-(-\eps_i-\eps_l+\eps_e-\eps_m)-\ii0^+}
\\
&+
\ii\sum_{me}
\frac{
V_{km\,ie}V_{le\,jm}
}
{-\Omega-(-\eps_j-\eps_l+\eps_e-\eps_m)-\ii0^+}
\end{split}
\label{eq:pp_GF2_D}
\end{equation}
which is the dynamic $D$ block of the GF2 pp BSE matrix.
Equations~\eqref{eq:pp_GF2_C}, \eqref{eq:pp_GF2_B}, and
\eqref{eq:pp_GF2_D} reproduce term by term the second-order pp kernel
derived independently from the anomalous GF2 self-energy.
The agreement includes the orbital permutations, relative signs, energy
denominators, and the static character of the coupling block.

This comparison also fixes the pair normalization.
In unrestricted orbital sums the pp BSE contains an
explicit factor $1/2$ because the ordered pairs $(p,q)$ and $(q,p)$ are
both present.
In the normalized antisymmetric basis with $p<q$, this double counting is
removed and no additional $1/2$ appears in the multiresolvent expression
\begin{equation}
 \ii\Phi_P^\pp
 =
 U_L^\pp R^\pp U_R^\pp
 \label{eq:pp_master_normalization_check}
\end{equation}
The successful recovery of all three independent pp-GF2 blocks provides
an independent verification of this normalization.

The two primitive channels therefore exhibit a completely parallel
frequency-integration structure
\begin{equation}
\begin{array}{ccc}
\eh & & \pp
\\[3pt]
A(\omega) & \longleftrightarrow & C(\Omega)
\\
B & \longleftrightarrow & B
\\
A^\dagger(-\omega) & \longleftrightarrow & D(-\Omega)
\end{array}
\label{eq:eh_pp_GF2_parallel}
\end{equation}
In both channels the diagonal blocks retain the bosonic frequency while
the off-diagonal coupling block becomes static after the second
frequency integration.
This provides one of the strongest independent validations of the
second-order multiresolvent construction.

%%%%%%%%%%%%%%%%%%%%%%%%%

\section{Channel consistency at second order}
\label{sec:channel_consistency}
%%%%%%%%%%%%%%%%%%%%%%%%%

The second-order construction can be checked in several independent ways.
The first check is crossing.
The complete transverse $\eh$ contribution must be obtained from
the direct $\eh$ reducible vertex through
Eq.~\eqref{eq:Phi_ehbar_crossing}.
This fixes both the orbital permutation and the frequency combination
$\nu'-\nu$.

The second check uses the $\pp$ convention.
The native pair propagator depends only on $\Omega_\pp$, and the exact
frequency transformation of Eq.~\eqref{eq:pp_frequency_map} fixes the
combination $\nu+\nu'+\omega$ in the direct $\eh$ convention.
One-sided frequency integration of this object must produce the same $\hhp$ and $\pph$
external-vertex poles obtained from the transverse $\eh$ contribution.

The third check is the recovery of the $\eh$ and $\pp$ GF2 kernels after
integrating over the remaining fermionic frequency, as shown in
Secs.~\ref{sec:second_contraction} and \ref{sec:pp_gf2_check}.
This check is independent of the initial channel-resolved derivation and
is sensitive to the signs, orbital-index permutations, and the $1/2$
factors in the pair channel.

Finally, the full second-order vertex must be independent of the channel
through which it is reconstructed.
Using the parquet relations,
\begin{align}
	\Gamma^{\eh,(2)}+\Phi^{\eh,(2)}
	&=
	\Phi^{\eh,(2)}+\Phi^{\teh,(2)}+\Phi^{\pp,(2)}
	\label{eq:F_second_eh}
	\\
	\Gamma^{\pp,(2)}+\Phi^{\pp,(2)}
	&=
	\Phi^{\eh,(2)}+\Phi^{\teh,(2)}+\Phi^{\pp,(2)}
	\label{eq:F_second_pp}
	\\
	F^{(2)}
	&=
	\Gamma^{\eh,(2)}+\Phi^{\eh,(2)}
	=
	\Gamma^{\pp,(2)}+\Phi^{\pp,(2)}
	\label{eq:channel_consistency_second}
\end{align}
The same channel organization reconstructs the direct $\eh$ Hedin
kernel through Eq.~\eqref{eq:Gamma_eh}.

%%%%%%%%%%%%%%%%%%%%%%%%%
%%%%%%%%%%%%%%%%%%%%%%%%%
\newpage

\section{Recipe through second order}
\label{sec:recipe_second}
%%%%%%%%%%%%%%%%%%%%%%%%%

The construction is now complete through second order and can be summarized
without invoking any third-order ingredient.

First, fix the three channel transfers
\begin{equation}
	\Omega_\eh=\omega
	\qquad
	\Omega_{\teh}=\nu'-\nu
	\qquad
	\Omega_\pp=\nu+\nu'+\omega
	\label{eq:recipe_second_transfers}
\end{equation}
and construct the two primitive reducible vertices
$\Phi^{\eh,(2)}$ and $\Phi_P^{\pp,(2)}$ in their natural channels.
The transverse contribution $\Phi^{\teh,(2)}$ is generated exclusively by
crossing.

Second, represent the two primitive vertices as
\begin{align}
	\ii\Phi^{\eh,(2)}
	&=
	U_L^{\eh,(1)}R^{\eh,(0)}U_R^{\eh,(1)}
	\label{eq:recipe_second_eh}
	\\
	\ii\Phi_P^{\pp,(2)}
	&=
	U_L^{\pp,(1)}R^{\pp,(0)}U_R^{\pp,(1)}
	\label{eq:recipe_second_pp}
\end{align}
The central resolvents describe neutral $N$-electron and pair
$(N\pm2)$-electron propagation, respectively.

Third, reconstruct the second-order channel-irreducible vertices from
\begin{align}
	\Gamma^{\eh,(2)}
	&=
	\Phi^{\teh,(2)}+\Phi^{\pp,(2)}
	\\
	\Gamma^{\pp,(2)}
	&=
	\Phi^{\eh,(2)}+\Phi^{\teh,(2)}
\end{align}
and perform a one-sided fermionic-frequency integration.
The resulting pair-space external vertices contain only
\begin{equation}
	N-1:\ \hhp
	\qquad\qquad
	N+1:\ \pph
	\label{eq:recipe_second_external_spaces}
\end{equation}
poles.
A second fermionic-frequency integration recovers the one-frequency GF2 BSE
kernels in both primitive channels.

Finally, the full second-order vertex can be reconstructed equivalently from
either primitive channel
\begin{equation}
	F^{(2)}
	=
	\Gamma^{\eh,(2)}+\Phi^{\eh,(2)}
	=
	\Gamma^{\pp,(2)}+\Phi^{\pp,(2)}
	\label{eq:recipe_second_F}
\end{equation}
This closes the second-order construction.
Everything that follows concerns exclusively the extension to third order.

\newpage

\section{Third-order multiresolvent extension}
\label{sec:multiresolvent_third}
%%%%%%%%%%%%%%%%%%%%%%%%%

The construction up to Sec.~\ref{sec:recipe_second} is complete through
second order.
We now extend the same multiresolvent ansatz by one perturbative order.
No new frequency architecture is introduced: the third-order terms arise
only from corrections to the two external vertices and to the central
resolvent.

For either primitive channel $r=\eh,\pp$, we expand
\begin{align}
	U_{L/R}^r
	&=
	U_{L/R}^{r,(1)}
	+
	U_{L/R}^{r,(2)}
	+
	\mathcal O(V^3)
	\label{eq:target_U_expansion}
	\\
	M^r
	&=
	M^{r,(0)}
	+
	M^{r,(1)}
	+
	\mathcal O(V^2)
	\label{eq:target_M_expansion}
\end{align}
The corresponding central resolvent becomes
\begin{align}
	R^r
	&=
	R^{r,(0)}
	+
	R^{r,(1)}
	+
	\mathcal O(V^2)
	\label{eq:target_R_expansion}
	\\
	R^{r,(1)}
	&=
	R^{r,(0)}M^{r,(1)}R^{r,(0)}
	\label{eq:target_R_first}
\end{align}
The complete third-order contribution required by the ansatz is therefore
\begin{equation}
	\begin{split}
	\ii\Phi^{r,(3)}
	={}&
	U_L^{r,(2)}R^{r,(0)}U_R^{r,(1)}
	+
	U_L^{r,(1)}R^{r,(0)}U_R^{r,(2)}
	\\
	&+
	U_L^{r,(1)}R^{r,(1)}U_R^{r,(1)}
	\end{split}
	\label{eq:target_third_order}
\end{equation}
The three terms correspond, respectively, to a second-order correction to
the left external vertex, a second-order correction to the right external
vertex, and a first-order correction to the central resolvent.
The remainder of the notes determines these three ingredients separately in
the $\eh$ and $\pp$ channels and verifies that they exhaust the complete
third-order reducible vertices.

\section{Direct eh vertex through third order}
\label{sec:eh_third}
%%%%%%%%%%%%%%%%%%%%%%%%%

It is useful to distinguish two closely related partially integrated objects.
The quantities $U^{++}_{ia,jb}$ and $U^{+-}_{ia,jb}$ derived in
Sec.~\ref{sec:half_contraction_second} are partially integrated pair-space
kernels and were introduced to test the reduction to the one-frequency
GF2 kernel.
The external vertices entering Eq.~\eqref{eq:target_eh} have instead one
general external index pair and one central neutral pair.
Both are obtained from the same three-frequency
$\Gamma^{\eh,(2)}$ and therefore possess the same $\hhp$ and $\pph$
pole spaces, but their external index grouping is different.

%%%%%%%%%%%%%%%%%%%%%%%%%
\subsection{Perturbative BSE expansion}
%%%%%%%%%%%%%%%%%%%%%%%%%

Define the reduced quantities
\begin{equation}
	\mathcal G^\eh
	=
	\ii\Gamma^\eh
	\qquad
	\mathcal P^\eh
	=
	\ii\Phi^\eh
	\label{eq:reduced_vertices}
\end{equation}
Before one-sided frequency integration, products of frequency-dependent vertices and pair
propagators are frequency convolutions.
We denote these convolutions by the symbol $\circ$.
The complete third-order direct $\eh$ reducible vertex is
\begin{equation}
	\begin{split}
	\mathcal P^{\eh,(3)}
	={}&
	\mathcal G^{\eh,(2)}
	\circ L_0
	\circ \mathcal G^{\eh,(1)}
	\\
	&+
	\mathcal G^{\eh,(1)}
	\circ L_0
	\circ \mathcal G^{\eh,(2)}
	\\
	&+
	\mathcal G^{\eh,(1)}
	\circ L_0
	\circ \mathcal G^{\eh,(1)}
	\circ L_0
	\circ \mathcal G^{\eh,(1)}
	\end{split}
	\label{eq:Phi_eh_third_convolution}
\end{equation}
No $\Gamma^{\eh,(3)}$ contribution appears in
$\Phi^{\eh,(3)}$ because an $\eh$ reducible diagram contains at
least one interaction on each side of the cut.
A third-order irreducible vertex would therefore first enter
$\Phi^\eh$ at fourth order.

The fully irreducible vertex also has no second- or third-order
contribution beyond the bare interaction,
\begin{equation}
	\Lambda
	=
	V
	+
	\mathcal O(V^4)
	\label{eq:Lambda_through_third}
\end{equation}
so that the second-order irreducible kernel entering
Eq.~\eqref{eq:Phi_eh_third_convolution} is exactly the object derived in
Eq.~\eqref{eq:Gamma_eh_second_explicit}.

The zeroth-order neutral resolvent $R^{\eh,(0)}$ and the first-order external vertices were fixed already in Sec.~\ref{sec:second_order_multiresolvent}.
The new ingredients required at third order are therefore the second-order external-vertex corrections and $R^{\eh,(1)}$.

%%%%%%%%%%%%%%%%%%%%%%%%%
\subsection{Second-order corrections to the external vertices}
%%%%%%%%%%%%%%%%%%%%%%%%%

The left forward external-vertex correction is obtained by integrating over one
fermionic frequency of the second-order irreducible kernel on its central neutral pair,
\begin{equation}
	(U_L^{\eh,(2),+})_{pr,ia}(\nu,\omega)
	=
	\frac{1}{2\pi}
	\int d\bar\nu
	\left[
	\ii\Gamma^{\eh,(2)}_{pa\,ri}
	(\nu,\bar\nu,\omega)
	\right]
	Q^+_{ia}(\bar\nu,\omega)
	\label{eq:UL_eh_second_def}
\end{equation}
The explicit result is
\begin{equation}
	\begin{split}
	(U_L^{\eh,(2),+})_{pr,ia}(\nu,\omega)
	={}&
	\sum_{jb}
	\frac{
	V_{pb\,ij}
	V_{aj\,rb}
	}
	{\nu-(\eps_i+\eps_j-\eps_b)-\ii0^+}
	\\
	&-
	\sum_{jb}
	\frac{
	V_{pj\,ib}
	V_{ab\,rj}
	}
	{\nu+\omega-(\eps_a+\eps_b-\eps_j)+\ii0^+}
	\\
	&-
	\frac{1}{2}
	\sum_{jk}
	\frac{
	V_{pa\,jk}
	V_{jk\,ri}
	}
	{\nu-(\eps_j+\eps_k-\eps_a)-\ii0^+}
	\\
	&+
	\frac{1}{2}
	\sum_{bc}
	\frac{
	V_{pa\,bc}
	V_{bc\,ri}
	}
	{\nu+\omega-(\eps_b+\eps_c-\eps_i)+\ii0^+}
	\end{split}
	\label{eq:UL_eh_second_explicit}
\end{equation}
The first two terms originate from the transverse $\eh$ channel
and the last two from the $\hh$ and $\pp$ branches.

The right forward external-vertex correction is
\begin{equation}
	(U_R^{\eh,(2),+})_{ia,sq}(\nu',\omega)
	=
	\frac{1}{2\pi}
	\int d\bar\nu
	Q^+_{ia}(\bar\nu,\omega)
	\left[
	\ii\Gamma^{\eh,(2)}_{sa\,qi}
	(\bar\nu,\nu',\omega)
	\right]
	\label{eq:UR_eh_second_def}
\end{equation}
which gives
\begin{equation}
	\begin{split}
	(U_R^{\eh,(2),+})_{ia,sq}(\nu',\omega)
	={}&
	\sum_{kc}
	\frac{
	V_{sk\,ic}
	V_{ac\,qk}
	}
	{\nu'-(\eps_i+\eps_k-\eps_c)-\ii0^+}
	\\
	&-
	\sum_{kc}
	\frac{
	V_{sc\,ik}
	V_{ak\,qc}
	}
	{\nu'+\omega-(\eps_a+\eps_c-\eps_k)+\ii0^+}
	\\
	&-
	\frac{1}{2}
	\sum_{kl}
	\frac{
	V_{sa\,kl}
	V_{kl\,qi}
	}
	{\nu'-(\eps_k+\eps_l-\eps_a)-\ii0^+}
	\\
	&+
	\frac{1}{2}
	\sum_{cd}
	\frac{
	V_{sa\,cd}
	V_{cd\,qi}
	}
	{\nu'+\omega-(\eps_c+\eps_d-\eps_i)+\ii0^+}
	\end{split}
	\label{eq:UR_eh_second_explicit}
\end{equation}
The backward external vertices follow from the same pair-reversal operation as
the pair-space external vertices.
The explicit formulas are not required for the third-order matching.

Equations~\eqref{eq:UL_eh_second_explicit} and
\eqref{eq:UR_eh_second_explicit} demonstrate that the external vertices of the
reducible direct $\eh$ object contain the same $\hhp$ and $\pph$
resolvents as the partially integrated Hedin kernel.
The difference lies only in the grouping of the external indices.

%%%%%%%%%%%%%%%%%%%%%%%%%
\subsection{First-order correction to the neutral resolvent}
%%%%%%%%%%%%%%%%%%%%%%%%%

The first-order neutral interaction matrix is
\begin{equation}
	M^{\eh,(1)}
	=
	\begin{pmatrix}
	A^{(1)}&B^{(1)}
	\\
	B^{(1)\dagger}&A^{(1)\dagger}
	\end{pmatrix}
	\label{eq:M_eh_first}
\end{equation}
with
\begin{align}
	A^{(1)}_{ia,jb}
	&=
	V_{aj\,ib}
	\label{eq:A_eh_first}
	\\
	B^{(1)}_{ia,bj}
	&=
	V_{ab\,ij}
	\label{eq:B_eh_first}
\end{align}
Using the general expansion of Eq.~\eqref{eq:target_R_first}, the first-order neutral-resolvent correction is
\begin{equation}
	R^{\eh,(1)}
	=
	R^{\eh,(0)}
	M^{\eh,(1)}
	R^{\eh,(0)}
	\label{eq:R_eh_first}
\end{equation}
No additional metric matrix appears in Eq.~\eqref{eq:R_eh_first} because the
metric is already included in the definition of the resolvent.

In the Tamm--Dancoff approximation (TDA), the corresponding central third-order
contribution is
\begin{equation}
	\begin{split}
	\left[
	\ii\Phi_{\rm cent}^{\eh,(3)}
	\right]_{pqrs}
	={}&
	\sum_{ia,jb}
	\frac{
	V_{pa\,ri}
	V_{aj\,ib}
	V_{qj\,sb}
	}
	{
	[\omega-(\eps_a-\eps_i)+\ii0^+]
	[\omega-(\eps_b-\eps_j)+\ii0^+]
	}
	\end{split}
	\label{eq:Phi_eh_third_central_tda}
\end{equation}
The full forward-backward expression is generated automatically by the
first-order correction to the neutral resolvent
\begin{equation}
	U_L^{\eh,(1)}
	R^{\eh,(1)}
	U_R^{\eh,(1)}
	=
	U_L^{\eh,(1)}
	R^{\eh,(0)}
	M^{\eh,(1)}
	R^{\eh,(0)}
	U_R^{\eh,(1)}
	\label{eq:Phi_eh_third_central_matrix}
\end{equation}

%%%%%%%%%%%%%%%%%%%%%%%%%
\subsection{Complete direct eh result through third order}
%%%%%%%%%%%%%%%%%%%%%%%%%

After the one-sided frequency integrations, the frequency convolutions in
Eq.~\eqref{eq:Phi_eh_third_convolution} reduce to ordinary matrix products.
The complete third-order result is
\begin{equation}
	\begin{split}
	\ii\Phi^{\eh,(3)}
	={}&
	U_L^{\eh,(2)}
	R^{\eh,(0)}
	U_R^{\eh,(1)}
	\\
	&+
	U_L^{\eh,(1)}
	R^{\eh,(0)}
	U_R^{\eh,(2)}
	\\
	&+
	U_L^{\eh,(1)}
	R^{\eh,(1)}
	U_R^{\eh,(1)}
	\end{split}
	\label{eq:Phi_eh_third_final}
\end{equation}
Together with Eq.~\eqref{eq:Phi_eh_second_resolvent}, this reproduces the
target representation of Eq.~\eqref{eq:target_eh} through third order.

The three terms in Eq.~\eqref{eq:Phi_eh_third_final} have a transparent
interpretation.
The first term corrects the left external vertex, the second corrects the right
external vertex, and the third corrects the central neutral resolvent.
These exhaust the third-order direct $\eh$ reducible diagrams.
Consequently, through third order all contributions are absorbed into
corrections to the two external vertices and the central neutral resolvent.

The explicit products contain genuine multiresolvent structures.
For example, a correction to the left external vertex contains terms of the form
\begin{equation}
	\frac{1}{\nu-E^-_{I^-}-\ii0^+}
	\frac{1}{\omega-\Delta\eps_{ia}+\ii0^+}
	\label{eq:multi_resolvent_left_minus}
\end{equation}
and
\begin{equation}
	\frac{1}{\nu+\omega-E^+_{I^+}+\ii0^+}
	\frac{1}{\omega-\Delta\eps_{ia}+\ii0^+}
	\label{eq:multi_resolvent_left_plus}
\end{equation}
while the central correction contains two neutral resolvents separated by
the first-order interaction matrix.

%%%%%%%%%%%%%%%%%%%%%%%%%
\section{pp vertex through third order}
\label{sec:pp_third}
%%%%%%%%%%%%%%%%%%%%%%%%%

The zeroth-order $\pp$ metric, pair matrix, central resolvent, and first-order
external vertices were established in
Sec.~\ref{sec:second_order_multiresolvent}.
The third-order extension requires the first-order correction to the pair
matrix and the second-order corrections to the external vertices.

The first-order pair interaction matrix is
\begin{equation}
	M^{\pp,(1)}
	=
	\begin{pmatrix}
	V_{ab\,cd}&V_{ab\,ij}
	\\
	V_{ij\,cd}&V_{ij\,kl}
	\end{pmatrix}
	\label{eq:M_pp_first}
\end{equation}
The upper-left block acts in the $N+2$ sector, the lower-right block in the
$N-2$ sector, and the off-diagonal blocks couple the two pair sectors.
Accordingly,
\begin{equation}
	R^\pp
	=
	R^{\pp,(0)}
	+
	R^{\pp,(1)}
	+
	\mathcal O(V^2)
	\label{eq:R_pp_expansion}
\end{equation}
with
\begin{equation}
	R^{\pp,(1)}
	=
	R^{\pp,(0)}
	M^{\pp,(1)}
	R^{\pp,(0)}
	\label{eq:R_pp_first}
\end{equation}

From Eq.~\eqref{eq:Gamma_pp_second}, the second-order corrections to the $\pp$ external vertices are generated by $\eh$ intermediate propagation.
One-sided frequency integration with the one-sided pp pair propagator produces
$N+1$ $\pph$ and $N-1$ $\hhp$ external-vertex poles.
For the two-particle central sector,
\begin{equation}
	U_{L,\ee}^{\pp,(2)}
	=
	\sum_{\ell}
	\sum_{I^+}
	\frac{
	\mathcal S^{+,\ell}_{I^+}
	}
	{z_\ell-E^+_{I^+}+\ii0^+}
	\label{eq:UL_pp_ee_second}
\end{equation}
while for the two-hole central sector,
\begin{equation}
	U_{L,\hh}^{\pp,(2)}
	=
	\sum_{\ell}
	\sum_{I^-}
	\frac{
	\mathcal S^{-,\ell}_{I^-}
	}
	{z_\ell-E^-_{I^-}-\ii0^+}
	\label{eq:UL_pp_hh_second}
\end{equation}
Here $z_\ell$ denotes the one-particle energy carried by either member of
the external pair and $\ell$ labels the two antisymmetrically related
attachments.
The corresponding right external vertices have the same pole spaces.
The detailed residues are not required for the third-order parquet
closure and are therefore left in channel-adapted form.

At third order,
\begin{equation}
	\begin{split}
	\ii\Phi_P^{\pp,(3)}
	={}&
	U_L^{\pp,(2)}
	R^{\pp,(0)}
	U_R^{\pp,(1)}
	\\
	&+
	U_L^{\pp,(1)}
	R^{\pp,(0)}
	U_R^{\pp,(2)}
	\\
	&+
	U_L^{\pp,(1)}
	R^{\pp,(1)}
	U_R^{\pp,(1)}
	\end{split}
	\label{eq:Phi_pp_third_final}
\end{equation}
The pure two-particle part of the central correction is the third-order
pair ladder
\begin{equation}
	\begin{split}
	\left[
	\ii\Phi_{P,\rm cent}^{\pp,(3)}
	\right]_{pqrs}^{\ee}
	={}&
	\sum_{\substack{a<b\\c<d}}
	\frac{
	V_{pq\,ab}
	V_{ab\,cd}
	V_{cd\,rs}
	}
	{
	[\Omega-(\eps_a+\eps_b)+\ii0^+]
	[\Omega-(\eps_c+\eps_d)+\ii0^+]
	}
	\end{split}
	\label{eq:Phi_pp_third_central_ee}
\end{equation}
The remaining $\ee\leftrightarrow\hh$ and $\hh\to\hh$ contributions are
generated by the other blocks of $M^{\pp,(1)}$.

Equation~\eqref{eq:Phi_pp_third_final} exhausts the third-order
$\pp$ reducible diagrams for the same reason as in the direct
$\eh$ channel.
The corrections to the external vertices contain the second-order pp-irreducible kernel and
the central correction contains one first-order $\pp$ random-phase-approximation (ppRPA) interaction.
Thus, through third order all $\pp$ reducible contributions are
absorbed into corrections to the two external vertices and the central pair
resolvent.

When expressed in the direct $\eh$ frequency convention, the
central pair transfer is
\begin{equation}
	\Omega_\pp
	=
	\nu+\nu'+\omega
	\label{eq:Omega_pp_repeated}
\end{equation}
which generates the characteristic anti-diagonal pair structure of the
three-frequency vertex.
\section{Consistency checks at third order}
\label{sec:checks_third}
%%%%%%%%%%%%%%%%%%%%%%%%%

The third-order extension is deliberately overdetermined.
Several independent checks can therefore be applied to the new third-order
terms, while the second-order checks have already been completed before
Sec.~\ref{sec:recipe_second}.

%%%%%%%%%%%%%%%%%%%%%%%%%
\subsection{Second frequency integration}
%%%%%%%%%%%%%%%%%%%%%%%%%

Integrating over the remaining fermionic frequencies in
Eq.~\eqref{eq:Phi_eh_third_final} gives
\begin{equation}
	\begin{split}
	\ii\widetilde\Phi^{\eh,(3)}
	={}&
	\widetilde\Gamma^{\eh,(2)}R^{\eh,(0)}\Gamma^{\eh,(1)}
	+
	\Gamma^{\eh,(1)}R^{\eh,(0)}\widetilde\Gamma^{\eh,(2)}
	\\
	&+
	\Gamma^{\eh,(1)}R^{\eh,(1)}\Gamma^{\eh,(1)}
	\end{split}
	\label{eq:third_one_frequency_check}
\end{equation}
which is precisely the third-order expansion of the one-frequency BSE with a
second-order dynamic kernel.

%%%%%%%%%%%%%%%%%%%%%%%%%
\subsection{Crossing}
%%%%%%%%%%%%%%%%%%%%%%%%%

The transverse $\eh$ primitive is never constructed
independently.
At every perturbative order,
\begin{equation}
	\Phi^{\teh,(n)}_{pqrs}(\nu,\nu',\omega)
	=
	-
	\Phi^{\eh,(n)}_{pqsr}
	\left(
	\nu,
	\nu+\omega,
	\nu'-\nu
	\right)
	\label{eq:crossing_order_n}
\end{equation}
This maps each left, right, and central third-order direct $\eh$
contribution into the corresponding transverse topology and preserves the
fermionic crossing symmetry order by order.

%%%%%%%%%%%%%%%%%%%%%%%%%
\subsection{Parquet reconstruction}
%%%%%%%%%%%%%%%%%%%%%%%%%

Because the fully irreducible vertex has no third-order correction,
\begin{align}
	F^{(3)}
	&=
	\Phi^{\eh,(3)}+\Phi^{\teh,(3)}+\Phi^{\pp,(3)}
	\label{eq:F_n_parquet}
	\\
	\Gamma^{\eh,(3)}
	&=
	\Phi^{\teh,(3)}+\Phi^{\pp,(3)}
	\label{eq:Gamma_eh_n}
	\\
	\Gamma^{\pp,(3)}
	&=
	\Phi^{\eh,(3)}+\Phi^{\teh,(3)}
	\label{eq:Gamma_pp_n}
	\\
	F^{(3)}
	&=
	\Gamma^{\eh,(3)}+\Phi^{\eh,(3)}
	=
	\Gamma^{\pp,(3)}+\Phi^{\pp,(3)}
	\label{eq:channel_consistency_n}
\end{align}
The direct $\eh$ Hedin kernel is therefore reconstructed from the
two primitive multiresolvent objects through
$\Gamma^\eh=\Lambda+\Phi^{\teh}+\Phi^\pp$.

%%%%%%%%%%%%%%%%%%%%%%%%%
\subsection{High-frequency structure}
%%%%%%%%%%%%%%%%%%%%%%%%%

The second-order corrections to the external vertices decay as the inverse of their fermionic frequency,
\begin{equation}
	U_L^{(2)}
	=
	\mathcal O(\nu^{-1})
	\qquad
	U_R^{(2)}
	=
	\mathcal O({\nu'}^{-1})
	\label{eq:side_high_frequency}
\end{equation}
When both fermionic frequencies are large, only the central
$\omega$-dependent contribution survives at third order.
The resulting hierarchy is therefore consistent with the standard
parquet asymptotic decomposition into a purely bosonic part and left and
right one-fermionic-frequency corrections.

The same argument applies in the pp channel with the central pair
frequency $\Omega_\pp$ held fixed.

%%%%%%%%%%%%%%%%%%%%%%%%%
\subsection{Completeness and absence of double counting}
%%%%%%%%%%%%%%%%%%%%%%%%%

Every third-order diagram reducible in the direct $\eh$ channel
belongs to one of three disjoint classes,
\begin{equation}
	\Gamma^{\eh,(2)}L_0\Gamma^{\eh,(1)}
	\qquad
	\Gamma^{\eh,(1)}L_0\Gamma^{\eh,(2)}
	\qquad
	\Gamma^{\eh,(1)}L_0\Gamma^{\eh,(1)}L_0\Gamma^{\eh,(1)}
	\label{eq:third_classes_eh}
\end{equation}
The first two contain a second-order block irreducible in the direct
$\eh$ channel, while the last is generated by repeated direct
$\eh$ scattering.
The corresponding statement holds in the $\pp$ channel.

The parquet classification assigns each two-particle reducible diagram to
one and only one reducibility channel.
The left external-vertex, right external-vertex, and central-resolvent classes above therefore do not
overlap.
Together with Eq.~\eqref{eq:Lambda_through_third}, this proves
perturbative completeness of the two primitive multiresolvent
constructions through third order.

%%%%%%%%%%%%%%%%%%%%%%%%%
\section{Recipe through third order}
\label{sec:recipe}
%%%%%%%%%%%%%%%%%%%%%%%%%

Starting from the complete second-order construction summarized in
Sec.~\ref{sec:recipe_second}, only three additional ingredients are required
for each primitive channel $r=\eh,\pp$.

First, obtain the second-order corrections to the left and right external
vertices from the corresponding second-order channel-irreducible kernel.
Their poles remain confined to the $(N-1)$ $\hhp$ and $(N+1)$ $\pph$
sectors.

Second, construct the first-order correction to the central matrix,
$M^{r,(1)}$, and hence
\begin{equation}
	R^{r,(1)}
	=
	R^{r,(0)}M^{r,(1)}R^{r,(0)}
	\label{eq:recipe_R_first}
\end{equation}
The central resolvent describes $N$-electron propagation for $r=\eh$ and
$(N\pm2)$-electron propagation for $r=\pp$.

Third, assemble the complete third-order reducible vertex
\begin{equation}
	\begin{split}
	\ii\Phi^{r,(3)}
	={}&
	U_L^{r,(2)}R^{r,(0)}U_R^{r,(1)}
	+
	U_L^{r,(1)}R^{r,(0)}U_R^{r,(2)}
	\\
	&+
	U_L^{r,(1)}R^{r,(1)}U_R^{r,(1)}
	\end{split}
	\label{eq:recipe_third_general}
\end{equation}
Generate $\Phi^{\teh,(3)}$ by crossing and reconstruct
\begin{align}
	\Gamma^{\eh,(3)}
	&=
	\Phi^{\teh,(3)}+\Phi^{\pp,(3)}
	\\
	F^{(3)}
	&=
	\Phi^{\eh,(3)}+\Phi^{\teh,(3)}+\Phi^{\pp,(3)}
\end{align}
Together with the second-order results, this yields the complete parquet
vertex and direct $\eh$ Hedin kernel through third order.
The two primitive channels retain the common particle-number architecture
\begin{align}
	\Phi^\eh
	&:\quad
	N\pm1\longrightarrow N\longrightarrow N\pm1
	\label{eq:sector_eh}
	\\
	\Phi^\pp
	&:\quad
	N\pm1\longrightarrow N\pm2\longrightarrow N\pm1
	\label{eq:sector_pp}
\end{align}
with common $(N\pm1)$ external-vertex pole spaces and channel-specific
central resolvents.

%%%%%%%%%%%%%%%%%%%%%%%%
\begin{thebibliography}{99}

\bibitem{Hedin1965}
L.~Hedin,
``New method for calculating the one-particle Green's function with application to the electron-gas problem,''
Phys. Rev. \textbf{139}, A796--A823 (1965)

\bibitem{DeDominicis_1964b}
C.~De~Dominicis and P.~C.~Martin,
``Stationary entropy principle and renormalization in normal and superfluid systems. II. Diagrammatic formulation,''
J. Math. Phys. \textbf{5}, 31--59 (1964)

\bibitem{Bickers_2004}
N.~E.~Bickers,
``Self-consistent many-body theory for condensed matter systems,''
in \textit{Theoretical Methods for Strongly Correlated Electrons},
edited by D.~S\'en\'echal, A.-M.~S.~Tremblay, and C.~Bourbonnais
(Springer, New York, 2004), pp.~237--296

\bibitem{MarieLoos2025}
A.~Marie and P.-F.~Loos,
``Parquet theory for molecular systems: Formalism and static kernel parquet approximation,''
J. Chem. Phys. \textbf{163}, 194115 (2025)

\bibitem{MarieAmmarLoos2024}
A.~Marie, A.~Ammar, and P.-F.~Loos,
``The $GW$ approximation: A quantum chemistry perspective,''
Adv. Quantum Chem. \textbf{90}, 157--184 (2024)

\end{thebibliography}
%%%%%%%%%%%%%%%%%%%%%%%%

\end{document}
